%%%%%%%%%%%%%%%%%%%%%%%%%%%%%%%%%%%%%%%%%%%%%%%%%%%%%%%%%%%%%%%%%%%
\chapter{Memory Circuits\index{memories:  first order digital systems (1-OS) having state autonomy due to the first loop closed over 0-OS}: One-Loop, First-Order Systems (1-OS)\index{1-OS: first order system -- memory}}
%%%%%%%%%%%%%%%%%%%%%%%%%%%%%%%%%%%%%%%%%%%%%%%%%%%%%%%%%%%%%%%%%%%%
\localtoc

In this chapter, we will introduce memory circuits that are obtained through a first loop that brings to the input of a circuit the effect of applying a signal to another input. This reverse connection allows the circuit to maintain a state according to a command. Thus, the memory function is obtained.

In this textbook, we will not cover every possible 1-OS circuit. We cover only those circuits that are used in building a processor. 

%%%%%%%%%%%%%%%%
\section{Latch}
%%%%%%%%%%%%%%%%

The simplest 1-OS circuit allows for storing the simplest event: the temporary transition of a digital signal from 1 to 0, or the temporary transition of some circuit from 0 to 1. A latch is a circuit that performs this function.

\placedrawing[H]{figures/03_01_latches.lp}{{\bf The elementary latches.}
{\small Using the loop, closed from the output to one input, elementary storage elements are built. {\bf a}. {\em AND
loop} provides a {\em reset-only} latch. {\bf b}. {\em OR loop}
provides the {\em set-only} version of a storage element. {\bf c}.
The heterogeneous elementary {\em set-reset} latch results
combining the {\em reset-only} latch with the {\em set-only}
latch. {\bf d}. The waveforms describe the behavior of the previous three latch circuits.}}{latches}


%************************************
\subsection{Closing the first loop}
%**************************************

Closing a feedback loop requires that an output of a logic circuit be connected to one of its inputs. In Figure \ref{latches}a, the output of an AND gate is connected to one of the inputs, and the pulse (Reset)' is applied to the other input. If the transition to zero lasts long enough so that the signal propagates through the circuit to the second input, then the output is fixed at the zero value, and the signal (Reset)' from the input can disappear. The transition event from the input is fixed to the output forever.

Now, let us take an OR gate and connect its output to one of the inputs. Then, we can apply a transition from 0 to 1. If the duration of the input signal allows propagation of its effect on the reaction loop, the output of the circuit will be permanently fixed to the value 1.

Latches\index{latch!elementary latch: based on an even inverting level combinational loop captures temporary events} allow us to capture temporary events indefinitely. We say that the two circuits memorize, one the temporary transition to 0, the other the temporary transition to 1. However, for the simple AND and OR loops, we cannot make these circuits "forget" what they have memorized. But, by combining the two circuits, we obtain a fundamental memory structure that can change its state according to the set command, S, to 1 or, according to the reset command, R', to 0. The set command is active on 1, while the reset command is active on 0.

\placedrawing[h]{figures/03_02_symmetric_latches.lp}{{\bf Symmetric elementary latches.}
{\small {\bf a.} Symmetric elementary {\em NAND latch} with
low-active commands S' and R'. {\bf b}. Symmetric elementary {\em
NOR latch} with high-active commands S and R.}}{symLatches}

Using both an AND and an OR, we obtain a circuit with two states between which it can switch according to asymmetrical commands, one active on zero and the other active on one. Often, it is preferable that both commands are of the same type. We can use the two forms of De Morgan's law to transform the circuit in Figure \ref{latches}c and we will obtain two symmetrical structures ordered with signals of the same type. Using the De Morgan law  A + B = (A'B')', the latch from Figure \ref{latches}c becomes the latch from Figure \ref{symLatches}a, and using the De Morgan law AB = (A'+B')', the latch from Figure \ref{latches}c becomes the latch from Figure \ref{symLatches}b.

%***************************
\subsection{Clocked Latch\index{latch!clocked latch: the Set/Reset latch transparent on the active level of clock}}
%****************************
In digital electronics, asynchronous SR latches can be problematic.  {\bf The first latch problem}: the inputs to indicate {\em how} the latch switches are the same as the inputs to
indicate {\em when} the latch switches. This immediate reactivity means they are highly susceptible to noise and glitches, especially when dealing with mechanical switches or noisy signals. A common issue is \textit{bouncing}, where the input rapidly fluctuates between high and low states due to imperfect physical contacts or noisy environments. This can cause an asynchronous SR latch to switch states unpredictably.

The first problem is solved in Figure \ref{ck_latch}a by conditioning the application of signals S'and R' to the latch in Figure \ref{symLatches}a by an additional signal that we will call the clock, CK. Thus, we will use the S and R inputs to specify how the circuit switches and the CK input to specify when the circuit switches. In this way, the decoupling of "{\em when}" from "{\em how}" is obtained. By synchronizing state changes with the clock, the system avoids unintended transitions caused by transient inputs or noise. 

For the clocked latch, two NAND gates are added (see Figure \ref{ck_latch}) to allow the S and R inputs to be conditioned by the clock.

\placedrawing[H]{figures/03_03_clocked_latches.lp}{{\bf Elementary clocked latch.}
{\small The transparent RS clocked latch is sensitive
(transparent) to the input signals during the active level of the
clock (the high level in this example). {\bf a}. The internal
structure. {\bf b}. The logic symbol.}}{ck_latch}

%***************************
\subsection{D Latch\index{latch!data latch (DL): a clocked latch with {\tt S = D} and {\tt R = D'}}}
%****************************
{\bf The second latch problem}: if we apply synchronously
S'=0 and R'=0 on the inputs of the NAND latch  (or S=1 and R=1 on the
inputs of the OR latch), i.e., the latch is commanded ``to switch in
both states simultaneously", then we cannot predict what is the
state of the latch after the ending of these two active signals.

For the second problem, we have a solution achieved through a restriction. We connect between the two inputs a NOT that will restrict the simultaneous application of setting and resetting the circuit (see Figure \ref{d_latch}a). 

\placedrawing[h]{figures/03_04_data_latch.lp}{{\bf The data latch.} {\small
Imposing the restriction $R=S'$ to an RS latch results in the {\bf D
latch} without non-predictable transitions ($R=S=1$ is not anymore
possible). {\bf a.} The structure. {\bf b.} The logic symbol. {\bf
c.} An improved version for the data latch internal
structure. {\bf d.} The optimal structure using an elementary multiplexer.}}{d_latch}

In the context of a latch, \textit{transparency} refers to the period during which the output follows the input directly. When the enable signal is active, any change at the input is immediately reflected at the output, making the latch behave like a simple wire. 

For the time interval in which the clock is active, CK = 1, it means that the latch is transparent to the control signal S and R. The latch becomes non-transparent in the time interval in which CK = 0. In the period in which the latch is transparent, the decoupling between "when" and "how" is not achieved, and the circuit switches are conditioned by any changes in the S and R inputs. A correct use of the latch assumes that during the transparency period, the S and R inputs are stable.



%%%%%%%%%%%%%%%%%%%%%%%%%%%%%%%%%%%%%%%%%%%%%%%%%%%
\section{Flip-Flops Based on the Primary-Secondary Principle\index{flip-flop!SR-FF: a flip-flop based on the Primary-Secondary principle}}
%%%%%%%%%%%%%%%%%%%%%%%%%%%%%%%%%%%%%%%%%%%%%%%%%%%

The clocked latch, in the SR version as well as in the D version (data latch), allows switching at any time during the transparency. This can result in an instability and/or unwanted changes in the output signal. Many applications require switching at a precise time, typically determined by the active, positive, or negative transition of the clock signal. Applying the {\em primary-secondary principle}\index{Primary-Secondary principle: the circuit configuration that ensures, by connecting two latches in series, the edge switching of the resulting storage element} will allow this behavior.

Figure \ref{m_s}a shows a structure in which transparency is blocked by two RS latches that have the clock applied in antiphase. On the active level of the clock, the positive one, the first latch, the primary latch, is transparent. This allows it to be switched according to the S and R inputs without affecting the output of the circuit because the second latch is not transparent. On the inactive level of the clock, the second latch copies the state of the first latch, which cannot be changed. In this way, the entire circuit switches as a consequence of the negative transition of the clock signal. So, the active edge of the clock applied to the primary-secondary structure is the negative one. In Figure \ref{m_s}b, the logical symbol of the structure in Figure \ref{m_s}a is represented.

\placedrawing[H]{figures/03_05_master_slave.lp}{{\bf The primary-secondary principle.}
{\small Serially connecting two RS latches, activated with
different levels of the clock signal, results in a non-transparent
storage element. {\bf a.} The structure of an RS primary-secondary
flip-flop, active on the falling edge of the clock signal. {\bf
b.} The logic symbol of the RS flip-flop triggered by the {\em negative edge} of clock. {\bf
c.} The logic symbol of the RS flip-flop triggered by the {\em positive edge} of clock.}}{m_s}

\placedrawing[H]{figures/03_06_d_flip_flop.lp}{{\bf The delay (D) flip-flop.}
{\small  Restricting the two inputs of an RS flip-flop to
$D=S=R'$, results an FF with predictable transitions. {\bf a.} The
structure. {\bf b.} The logic symbol. {\bf c.} The wave forms
proving the delay effect of the D flip-flop.}}{d_ff}

The second latch problem also propagates at the level of the primary-secondary structure. The solution we can come up with at this level is the one that introduces the limitation by connecting the inputs through an inverter circuit (see Figure \ref{d_ff}a). The result is the \textit{D flip-flop} circuit \index{flip-flop!D-FF: delay flip-flop, a RS-FF with {\tt S = D} and {\tt R = D'}}, D-FF (delay flip-flop).  A D flip-flop is edge-triggered, which means that it captures the input value only at the moment of a clock edge (either rising or falling). Once the clock edge passes, further changes in the input do not affect the output until the next clock edge.  The name is justified by the waveforms represented in Figure \ref{d_ff}c, where the circuit output follows the evolution of the input with a delay equal to the period of the clock signal.

%%%%%%%%%%%%%%%%%%%%%%%%%%%%%%%%%%%%%%%%%%%%%%
\section{Serial-parallel extension: Register\index{register: a $n$-bit storage device triggered by the active edge of the clock}}
%%%%%%%%%%%%%%%%%%%%%%%%%%%%%%%%%%%%%%%%%%%%%%

%*********************
\subsection{Structure}
%*********************

The serial-parallel extension in first-order systems (1-OS) is represented by the register. By connecting in parallel $n$ serially extended structures (the delay flip-flop type primary-secondary circuit) we obtain a {\bf {\em register}}. By applying the same clock signal on each D-FF (see Figure \ref{reg}a), an $n$-bit input vector is stored simultaneously. 

\placedrawing[h]{figures/03_07_register.lp}{{\bf The $n$-bit register.} {\small
{\bf a.} The structure: $n$ D-FFs connected in parallel.
{\bf b.} The logic symbol.} }{reg}

Figure \ref{register} shows the effect of the transfer through a register.

\placedrawing[h]{figures/03_08_register_operation.lp}{{\bf Register operation.} {\small At each active
edge of the clock (in this example, it is the positive edge), the register's
output takes the value applied on its inputs if {\tt reset = 0} and
{\tt enable = 1}.}}{register}

The time behavior is characterized by two time intervals:
\begin{description}
    \item[minimal set-up time]: $t_{SU_{min}}$ is the minimum time interval in which the input must remain stable before the active edge of the clock so that the transition of the output corresponds to the input
    \item[maximal propagation time]: $t_p$ is the propagation time of the output related to the active edge of the clock.
\end{description}

\ifthenelse{\boolean{chiselversion}}{
\includechiselcode[caption={Parameterized Register (Chisel: Register.scala)}, label={lst:chisel_register}, firstline=8, firstnumber=8, lastline=27]{\chiselSourceDir/memory/Register.scala}

Listing \ref{lst:chisel_register} shows parameterized Chisel code that implements an \textit{n}-bit register with an enable and reset signal. In Chisel, registers are created using the built-in \texttt{Reg} and \texttt{RegInit} classes, which define stateful elements that store values across clock cycles. The \texttt{RegInit(0.U(n.W))} statement initializes a register of width \textit{n} to zero. This ensures that the register has a defined state at power-up or reset.

One key feature of Chisel's register implementation is the \textit{inferred clock signal}. Unlike Verilog, where explicit clock signals must be defined, Chisel modules automatically use the system clock unless specified otherwise. This means that all registers declared within a module implicitly update on the \textit{rising edge} of the global clock. Consequently, the \texttt{reg} variable in this implementation updates its stored value only at the active clock edge when \texttt{io.enable} is high.

The use of \texttt{when(io.reset) { reg := 0.U }} ensures that the register resets to zero when \texttt{io.reset} is asserted. Otherwise, when \texttt{io.enable} is high, the register updates to store the value of \texttt{io.d}. The reset is synchronous, meaning the reset only takes effect on the active clock edge. Unlike asynchronous resets in Verilog, which can immediately change the register’s value when the reset signal is asserted, Chisel’s \texttt{RegInit} ensures that the reset occurs in a clocked manner. In Chisel, to make the reset asynchronous, you need to use reset types explicitly. For example, \texttt{val reg = withReset(io.reset)(RegInit(0.U(n.W)))
}

} %end Chisel

%**************************
\subsection{Applications}
%**************************
In this Section, we describe applications of the register in digital systems.

\subsubsection{Storing}
Adding an {\tt enable} input allows us to determine when (i.e., in what
clock cycle) the input is loaded into a register. If {\tt enable =
0}, the registers {\em store} the data loaded in the last clock
cycle when the condition {\tt enable = 1} was fulfilled. This
means that we can keep the content stored in the register as
 long as needed.

\subsubsection{Buffering}
A buffer is a circuit element that is used to strengthen a signal, isolate circuit components, or introduce a delay. It is essentially a single-input, single-output gate that outputs the same value as the input but provides a higher current drive and higher signal integrity. Registers can be used to {\em buffer} two distinct blocks so that some behaviors are not
transmitted through the register. For example, in Figure
\ref{register} the transitions from {\tt c} to {\tt d} and from
{\tt d} to {\tt e} at the input of the register are not
transmitted to the output.

\subsubsection{Synchronizing}
Digital signals are often generated asynchronously relative to system input, meaning they arrive at unpredictable times. However, for proper system operation, these signals must be received in a well-controlled and time-aligned manner. In such cases, we say the signals are applied \textit{asynchronously} but must be received \textit{synchronously}. 

For example, in Figure \ref{register}, the input to the register changes at arbitrary times with respect to the active edge of the clock. However, the output of the register transitions only with a constant delay after the positive clock edge. This behavior ensures that the input signals are synchronized with the register output, providing a consistent timing relationship. We refer to this process as \textit{time tempering}, as it transforms uncoordinated signal changes into a predictable and controlled timing structure.

\subsubsection{Delaying}
When a value is applied to the input of a register during clock cycle {\tt n}, it appears at the output of the register in clock cycle {\tt n+1}. In other words, the register introduces a delay of one clock cycle between its input and output. This behavior ensures that the output always reflects the input of the previous clock cycle. Figure \ref{register} illustrates this concept, showing how a value applied to the input of the register in one clock cycle is propagated to its output in the subsequent clock cycle.


\subsubsection{Looping}
Designing a digital system involves creating various types of connections between components. One particularly significant type is a connection in which certain outputs are fed back to specific inputs within a digital subsystem. These connections are known as {\bf loops}. Registers play an important role in the formation of controllable loops within complex systems, ensuring proper sequencing and synchronization of data flow.

%*************************
\subsubsection{Pipelining}
%**************************
\begin{exemplu}
The previously exemplified {\tt threeAdder.sv} module performs the addition of three numbers in twice the time of the sum of two numbers. We can reduce the execution time through a pipeline solution that uses two registers. The solution is presented in the following module:
{\em
\includeverilogcode[caption={The module 'threeAdder' has 3 4-bit inputs and one 4-bit output. The circuit adds modulo 16 three numbers; do not provide carry output.   File: A09\_pipelinedThreeAdder.sv  (SystemVerilog)}, 
                    label={lst:A09_pipelinedThreeAdder.sv}, 
                    firstline=1, 
                    firstnumber=1, 
                    lastline=60]
                    {\verilogSourceDir/A09_pipelinedThreeAdder.sv}
}  
Two modules of {\tt adder} defined Example \ref{adder}
are instantiated as {\tt inAdder} and {\tt outAdder}, they are
interconnected using the pipeline  register {\tt pipeReg}. The register {\tt subcReg} is used to synchronize the {\tt in2} input with the pipelined result provided by {\tt inAdder}.

$\diamond$
\end{exemplu}

%%%%%%%%%%%%%%%%%%%%%%%%%%%%%%%%%%%%%%%%%%%%%%%%%%%
\section{Parallel extension: Random-Access Memory\index{RAM (Random-Access Memory)!asynchronous  RAM}}
%%%%%%%%%%%%%%%%%%%%%%%%%%%%%%%%%%%%%%%%%%%%%%%%%%%

The parallel extension in first-order systems (1-OS) is represented by random access memory (RAM).

%******************************
\subsection{Generic structure}
%********************************

The generic structure of a memory with $2^p$ one-bit locations is represented in Figure \ref{ram}.

\placedrawing[h]{figures/03_09_ram.lp}{{\bf The principle of the random
access memory (RAM).} {\small The clock is distributed by a DMUX
to one of $m=2^{p}$ DLs, and the data is selected by a MUX from
one of the $m$ DLs. Both, DMUX and MUX use as selection code a
$p$-bit address. The one-bit data DIN can be stored in the clocked
DL.}}{ram}

The waveforms that define the operation of the structure in Figure \ref{ram} are represented in Figure \ref{wave} where the main restrictions are represented by:
\begin{description}
    \item[$t_{ASU}$]: address set-up time represents the time interval in which the address must be stable before the activation of the WE signal to allow the decoder in the DMUX to do its work.
    \item[$t_{DSU}$]: data set-up time represents the time interval in which DIN must be stable before de-activating the WE signal to allow the correct closing of the reaction loop in the selected latch.
    \item[$t_{DH}$]: data hold represents the time interval in which DIN must be stable after deactivating the WE signal to allow the correct closing of the reaction loop in the selected latch.

\placedrawing[H]{figures/03_10_async_ram.lp}{{\bf Read and write cycles for an
asynchronous RAM.} {\small Reading is a combinational process of
selecting. The access time, $t_{ACC}$, is given by the propagation
through a big MUX. The write enable signal must be strictly
included in the time interval when the address is stable (see
$t_{ASU}$ and $t_{AH}$). Data must be stable related to the
positive transition of $WE'$ (see $t_{DSU}$ and $t_{DH}$).}}{wave}

    
    \item[$t_{W}$]: width of the write enable signal which ensures correct writing in the selected latch
    \item[$t_{AH}$]: address hold time is the time interval in which the address must be maintained after the WE signal is provided
    \item[$t_{ACC}$]: access time is the time interval after which the content of the selected cell is accessible at the output
\end{description}
All the times previously listed are minimum values.

How can you build a RAM memory that stores words of $m$ bits? The generic solution is presented in Figure \ref{bigword} where for each bit a column containing latches $2^p$ is defined.

In Figure \ref{bigword}, $DCD_p$ is shared between the DMUX and the MUX in Figure \ref{ram}. The ANDs associated with the demultiplexing function are selected with the WE signal. The AND-OR structure is repeated $m$ times, once in each $COLUMN_i$ column.

\placedrawing[H]{figures/03_11_async_mbit_ram.lp}{{\bf The asynchronous $m$-bit word
RAM.} {\small Expanding the number of bits per word means to
connect in parallel one-bit word memories which share the same
decoder. Each COLUMN contains storing latches and  AND-OR
circuits for one bit.}}{bigword}\textcolor[rgb]{0.00,0.00,0.00}{}

%******************************************
\subsection{Synchronous RAM}\index{RAM (Random-Access Memory)!synchronous  RAM}\label{syncMem}
%*******************************************
At very high speeds, on the order of GHz, the time restrictions illustrated in Figure \ref{wave} are very difficult to fulfill in the asynchronous version described above. To increase the degree of controllability of our projects, synchronous RAM memories are used. In Figure \ref{swave}, the behavior of a synchronous memory (SRAM) is illustrated in the sense that all the time intervals that characterize the behavior of the memory are related to the active edge of a clock signal. In this case, the designer of a system that includes synchronous memory can more rigorously control the time behavior of the system.

For the memories used by the system designer, {\em memory libraries} are provided by specialized companies in the form of IPs (intellectual property), guaranteeing the correctness of the code used in the description.

\placedrawing[h]{figures/03_12_sync_ram.lp}{{\bf Read and write cycles for  Synchronous RAM (SRAM).}
{\small For the {\em flow-through} version of a $SRAM$ the time
behavior is similar to a register. The set-up and hold time are
defined related to the active edge of clock for all the input
connections: {\em data}, {\em write-enable}, and {\em address}.
The data output is also related to the same edge.}}{swave}

\includeverilogcode[caption={4096 32-bit words synchronous RAM; asynchronous read.   File: A09\_syncRAM.sv  (SystemVerilog)}, 
                    label={lst:A09_syncRAM.sv}, 
                    firstline=1, 
                    firstnumber=1, 
                    lastline=60]
                    {\verilogSourceDir/A09_syncRAM.sv}


\ifthenelse{\boolean{chiselversion}}{
In Chisel, memory can be implemented using built-in constructs such as \texttt{Mem} and \texttt{SyncReadMem}. For SRAM-like behavior, \texttt{SyncReadMem} is preferred since it introduces synchronous reads.

\includechiselcode[caption={Parameterized SRAM (Chisel: SRAM.scala)}, label={lst:chisel_sram}, firstline=9, firstnumber=9, lastline=999]{\chiselSourceDir/memory/SRAM.scala}

Listing \ref{lst:chisel_sram} shows a Chisel module that defines a synchronous SRAM (Static Random Access Memory) model, which provides controlled read and write operations using a clocked memory structure. The design consists of an addressable memory space, data input and output ports, and a write-enable control signal.

The \texttt{addr} input determines the memory location being accessed, and its bit width is computed using \texttt{log2Ceil(depth).W}, ensuring that it has enough bits to address all locations in the memory array. The \texttt{dataIn} and \texttt{dataOut} signals represent the data being written to or read from memory, while \texttt{writeEnable} controls whether a write operation is performed.

The core memory structure is implemented using \texttt{SyncReadMem}, a Chisel construct that models a synchronous memory block. This ensures that memory reads and writes occur only on clock edges, mimicking the behavior of actual SRAM in hardware.

Write operations are controlled using the \texttt{when} condition, which allows data to be written into memory only when \texttt{writeEnable} is asserted. The \texttt{mem.write(io.addr, io.dataIn)} statement stores the value of \texttt{dataIn} at the specified address.

For read operations, the module ensures that data is only retrieved when no write is occurring. The statement \texttt{io.dataOut := mem.read(io.addr, !io.writeEnable)} prevents read-after-write hazards by ensuring that reads do not interfere with ongoing write operations.

This SRAM model is particularly useful for FPGA-based Block RAM (BRAM) inference and can also be mapped to ASIC memory blocks during synthesis. However, while \texttt{SyncReadMem} may be inferred as SRAM in FPGA synthesis, ASIC tools require specific steps to ensure correct inference.

For ASIC synthesis, memory inference depends on the synthesis tool recognizing \texttt{SyncReadMem} as SRAM. However, inference is not always guaranteed. The best approach for large ASIC memories is to directly instantiate an SRAM macro provided by the foundry using \texttt{class HardMacroSRAM extends BlackBox }.
} %End Chisel


%**********************************************************
\subsection{Synchronous pipelined RAM}\index{RAM (Random-Access Memory)!synchronous pipelined RAM}\label{sybcPipeMem}
%**********************************************************
The time $t_{ACC}$ offered by a RAM can sometimes be too high for the speed requirements of the system. To increase the clock frequency, a register is used in the memory output. Thus, the data are obtained at the new output very quickly: $t_{ACC}$ is the propagation time through the register. But there is a price: the {\bf {\em latency}} of one clock cycle (we remember that the register consists of {\em delay} flip-flops).

A skilled designer will almost always be able to hide the effect of latency and take advantage of the increase in system frequency obtained through the pipeline register from the RAM output.

\includeverilogcode[caption={4096 32-bit words synchronous RAM; asynchronous read.   File: A09\_syncPipeRAM.sv  (SystemVerilog)}, 
                    label={lst:A09_syncPipeRAM.sv}, 
                    firstline=1, 
                    firstnumber=1, 
                    lastline=60]
                    {\verilogSourceDir/A09_syncPipeRAM.sv}

\ifthenelse{\boolean{chiselversion}}{

\includechiselcode[caption={Parameterized Pipelined SRAM (Chisel: PipelinedSRAM.scala)}, label={lst:chisel_sram_pipelined}, firstline=9, firstnumber=9, lastline=999]{\chiselSourceDir/memory/SRAM.scala}

Listing \ref{lst:chisel_sram_pipelined} implements a pipelined synchronous SRAM, where memory read operations have a one-cycle delay before the output is available. This delay is achieved using the \texttt{RegNext} construct, which ensures that the read data is stored in a register before being passed to the output.

The module defines an input-output interface consisting of several key signals. The \texttt{addr} signal specifies the memory address for both read and write operations. The \texttt{dataIn} signal provides the data to be stored in memory, while \texttt{dataOut} represents the registered output data, ensuring a pipelined read operation. The \texttt{writeEnable} signal determines whether a write occurs at the given address.

The memory itself is instantiated using \texttt{SyncReadMem}, a Chisel construct that models a clocked RAM where both read and write operations are synchronized with the clock. A write operation takes place when \texttt{writeEnable} is high, allowing the input data to be written to the specified memory address using \texttt{mem.write(io.addr, io.dataIn)}. Read operations, however, do not directly output data from memory. Instead, the read result is passed through a register before being assigned to the output.

A key aspect of this implementation is the use of \texttt{RegNext}. This introduces a pipeline stage that stores the read data in a register before it is sent to the output. Pipelining read data in this way reduces critical path delays by ensuring that the memory read operation occurs in one clock cycle while allowing the output to be registered in the next cycle. This improves timing closure in high-speed designs, preventing potential setup time violations. Additionally, the registered read output ensures stability by mitigating glitches that could arise from rapid address changes. 

The registered read data is assigned to the output using \texttt{io.dataOut := readReg}. This guarantees that the memory data is presented in a stable and predictable manner. The registered output ensures that data remains valid across clock cycles, making this design beneficial for both ASIC and FPGA implementations that require synchronous memory operations with controlled timing.

} %End Chisel



%**********************************************************
\subsection{{\bf {\em Register file}}\index{register file: a small, fast synchronous multi-port RAM}}\label{registerFile} 
%**********************************************************

A register file is a battery of registers from which, in each clock cycle, any two registers can be accessed for reading and any register for writing. The implementation of such a circuit is done with a synchronous memory with three ports: two for reading and one for writing (see Figure \ref{file_reg}), where:

\placedrawing[h]{figures/03_13_reg_file.lp}{{\bf Register file.} {\small In this
example it contains $2^{n}$ $m$-bit registers. In each clock cycle
any two registers can be read and writing can be performed in
anyone. }}{file_reg}

\begin{description}
    \item{{\tt leftAddr[n-1:0]}}: is the address used to select the data for the {\tt leftOut[m-1:0]} output
    \item{{\tt rightAddr[n-1:0]}}: is the address used to select the data for the {\tt rightOut[m-1:0]} output
    \item{{\tt destAddr[n-1:0]}}: is the address used to select the location where the value applied on input {\tt in[m-1:0]} is stored at the positive edge of the {\tt clock} signal
    \item{{\tt writeEnable}}: is the write enable command.
\end{description}
The internal structure of a file register is characterized by two multiplexers that ensure access to any pair of stored values.

\includeverilogcode[caption={Register file.   File: A09\_regFile.sv  (SystemVerilog)}, 
                    label={lst:A09_regFile.sv}, 
                    firstline=1, 
                    firstnumber=1, 
                    lastline=30]
                    {\verilogSourceDir/A09_regFile.sv}

\begin{verisum}{\em :

\hspace{-0.8cm} \rule[5mm]{16cm}{0.25mm} \vspace{-0.7cm}

\begin{description}
    \item[logic] [m-1:0] mem[0:n-1]: defines a block of memory of $n$ $m$-bit words
    \item[always\_ff]: defines a block which describes the behavior of register-like sequential circuits; it is followed always at least by {\tt {\bf @(posedge} clock)} or {\tt {\bf @(negedge} clock)} specifying the clock signal and its active edge. The non-blocking assignment, {\tt <=}, is used to assure the sequential behavior. While {\tt always\_comb} describes the behavior of a combinational circuit, {\tt always\_ff} ({\tt ff} comes from flip-flop) describes the behavior of the sequential, clock triggered, circuit called {\tt register}. Now {\tt logic} represents a register instead of an output of a combinational circuit like in the case of {\tt always\_comb}.
\end{description}
\rule[5mm]{16cm}{0.25mm} }
\end{verisum}

\ifthenelse{\boolean{chiselversion}}{
\includechiselcode[caption={Register File (Chisel: RegFile2R1WVec.scala)}, label={lst:chisel_regfile}, firstline=24, firstnumber=24, lastline=999]{\chiselSourceDir/memory/RegFile2R1WVec.scala}

Listing \ref{lst:chisel_regfile} implements a register-based register file. In Chisel, a register file can be implemented using either a \texttt{Vec} of \texttt{Reg} or an SRAM-based memory structure such as \texttt{SyncReadMem}. The choice between these two implementations depends on multiple factors, including the size of the register file, read and write access patterns, and the target hardware architecture.

A \texttt{Vec} of \texttt{Reg} is a collection of registers, each implemented as a distinct flip-flop. This approach is particularly suitable for small register files, such as general-purpose registers in a CPU core or pipeline buffers. Since each register is directly accessible, read operations are purely combinational, resulting in zero-cycle latency. This makes \texttt{Vec} of \texttt{Reg} ideal when low-latency register access is a requirement. Additionally, this structure allows multiple concurrent reads and writes, making it beneficial for designs where multiple functional units require simultaneous register access. However, as the number of registers increases, the area and power consumption grow significantly due to the need for dedicated flip-flops for every bit in the register file. This makes \texttt{Vec} of \texttt{Reg} inefficient for large-scale register files.

\texttt{SyncReadMem} provides a more scalable approach for implementing register files when the number of registers is large. It models a synchronous memory that stores register values in a memory array instead of individual flip-flops, significantly reducing area and power consumption. This is particularly advantageous when the register file must hold a large number of entries, such as in vector processors, GPU register banks, or out-of-order execution buffers. Since \texttt{SyncReadMem} models SRAM behavior, it typically has a one-cycle read latency, meaning that register values are not immediately available in the same cycle as the read request. This trade-off makes it less suitable for applications requiring immediate read results but is acceptable in architectures that can tolerate pipelined reads. Another limitation of \texttt{SyncReadMem} is that it typically supports only a single read and a single write port per clock cycle unless explicitly instantiated as multi-port memory, which can introduce complexity.

Additional designs showing SRAM register files can be found in Section \ref{subsec:regfile_2R1W_SRAM} and \ref{subsec:regfile_MT_2R1W_SRAM}.

}% End Chisel




%****************************************************
\subsection{Content-Addressable Memory\index{CAM (Content-Addressable Memory)}}\label{lb:cam}
%*****************************************************
A normal way to ``question" a memory circuit is to ask for:
\begin{quote}
{\bf Q1}: {\em what is the value of the property A of the object B}
\end{quote}
For example: {\em How old is George?} The {\em age} is the
property, and the {\em object} is George. The first step to design
an appropriate device to be questioned as previously is
exemplified is to define a circuit able to answer the question:
\begin{quote}
{\bf Q2}: {\em where is the object B?}
\end{quote}
with two possible answers:
\begin{enumerate}
    \item {\em the object B is not in the searched space}
    \item {\em the object B is stored in the cell indexed by X}.
\end{enumerate}
The circuit for answering Q2-type questions is called {\em Content
Addressable Memory}, {\em CAM}. (About Question Q1 in
the next subsection.)

The basic cell of a CAM consists of:
\begin{itemize}
    \item the storage elements for binary objects
    \item the ``questioning" circuits for searching for the value
    applied to the input of the cell.
\end{itemize}

\placedrawing[H]{figures/06_05_content_addressable_cell.lp}{{\bf Content Addressable Cell.}
{\small {\bf a.} The structure: data latches whose content is
compared against the input data using 4 XORs and one NAND. Write
is performed applying the clock with stable data input. {\bf b.}
The logic symbol.}}{cam_cell}

In Figure \ref{cam_cell} there are 4 D latches as storage elements
and four XORs connected to a 4-input NAND used as comparator. The
cell has two functions:

\begin{itemize}
    \item {\bf to store}: the active level of the clock modify the content
    of the cell storing  the 4-bit input data into the four D
    latches
    \item {\bf to search}: the input data are continuously compared
    with the cell content that generates the signal $AO'=0$ if
    the input matches the content.
\end{itemize}

\placedrawing[H]{figures/06_06_content_addressable_memory.lp}{{\bf The Content Addressable Memory
({\em CAM}).} {\small {\bf a.} A 4-word CAM is built using 4
content addressable cells, a demultiplexer to distribute the write
enable (WE) signal, and a $NAND_{4}$ to generate the match signal
(M). {\bf b.} The logic symbol.}}{cam}

The cell is {\em written} as a $m$-bit latch and is continuously
{\em interrogated} using a combinational circuit as a comparator.
The resulting circuit is an 1-OS because  results serially
connecting a memory with a combinational,
no-loop circuit. No additional loop is involved.

An $n$-word CAM contains $n$ CAM cells and some additional
combinational circuits to distribute the clock to the selected
cell and to generate the global signal M, activated to
signal a successful match between the input value and the contents of one or
more cells. In Figure \ref{cam}a a 4-word of 4 bits each
is represented. The write enable signal, $WE$, is demultiplexed as
clock to the appropriate cell, according to the address coupled by
$A_{1}A_{0}$. The 4-input NAND generate the signal $M$. If, at
least one address output, $AO_{i}'$ is zero, indicating match in
the corresponding cell, then $M = 1$ indicating a successful
search.


The {\em input address} $A_{p-1}, \ldots, A_{0}$ is binary codded
on $p = log_{2} n$ bits. The {\em output address} $AO_{n-1},
\ldots, AO_{0}$ is an {\em unary code} indicating the place or the
places where the data input $D_{m-1}, \ldots, D_{0}$ matches the
content of the cell. The output address must be unary codded
because there is the possibility of match in more than one cell.

Figure \ref{cam}b represents the logic symbol for a CAM with $n$
$m$-bit words. The input $WE$ indicate the function performed by
CAM. Be very careful with the set-up time and hold time of data
related to the $WE$ signal!

The CAM device is used to locate an object (to answer the question
{\bf Q2}). Dealing with the properties of an object (answering
{\bf Q1}-type questions) means to use o more complex devices which
{\em associate} one or more properties to an object. Thus, the
{\em associative memory} will be introduced adding some circuits
to CAM.


\subsection{Associative Memory\index{associative memory}}\label{assocMemory}

A partially used RAM can be an associative memory, but a very
inefficient one. Indeed, let be a RAM addressed by $A_{n-1} \ldots
A_{0}$ containing 2-field words $\{V, \; D_{m-1} \ldots D_{0}\}$.
The {\em objects} are codded using the address, the {\em values}
of the unique property $P$ are codded by the data field $D_{m-1}
\ldots D_{0}$. The one-bit field $V$ is used as a validation flag.
If $V=1$ in a certain location, then there is a match between the
object designated by the corresponding address and the value of
property $P$ designated by the associated data field.

\begin{exemplu}
Let be the 1Mword RAM addressed by $A_{19} \ldots A_{0}$
containing 2-field 17-bit words $\{V, \; D_{15} \ldots D_{0}\}$.
The set of objects, $OBJ$, are codded using 20-bit words, the
property $P$ associated to $OBJ$ is codded using 16-bit words. If
$$RAM[11110000111100001111] = 1\_0011001111110000$$
$$RAM[11110000111100001010] = 0\_0011001111110000$$
then:
\begin{itemize}
    \item for the object $11110000111100001111$ the property $P$
    is defined ($V=1$) and has the value $0011001111110000$
    \item for the object $11110000111100001010$ the property $P$
    is not defined ($V=0$) and the data field is meaningless.
\end{itemize}
Now, let us consider the 20-bit address codes four-letter names
using for each letter a 5-bit code. How many locations in this
memory will contain the field $V$ instantiated to 1?
Unfortunately, only extremely few of them, because:
\begin{itemize}
    \item only 24 from 32 binary configurations of 5 bits will be
    used to code the 24 letters of Latin alphabet ($24^{4} <
    2^{20}$)
    \item but more important: how many different name expressed by 4
    letters can be involved in a real application? Usually no more
    than few hundred, meaning almost nothing related to $2^{20}$.
\end{itemize}
 $\diamond$
\end{exemplu}

The previous example teaches us that a RAM used as associative
memory is a very inefficient solution. In real applications are
used names codded very inefficiently:
$$number\_of\_possible\_names >>> number\_of\_actual\_names.$$
In fact, the natural memory function means almost the same: to
remember about something immersed in a huge set of possibilities.

One way to implement an efficient associative memory is to take a
CAM and to use it as a {\em programmable decoder} for a RAM. The
(extremely) limited subset of the actual objects are stored into a
CAM, and the address outputs of the CAM are used instead of the
output of a combinational decoder to select the accessed location
of a RAM containing the value of the property $P$. In Figure
\ref{am} this version of an associative memory is presented.
$CAM_{m \times n}$ is usually dimensioned with $2^{m} >>> n$
working as a decoder programmed to decode {\bf any} very small
subset of $n$ addresses expressed by $m$ bits.

\placedrawing[h]{figures/06_07_associative_memory.lp}{{\bf An associative memory ({\em
AM}).} {\small The structure of an AM can be seen as a RAM with a
programmable decoder implemented with a CAM. The decoder is
programmed loading CAM with the considered addresses.}}{am}

Here are the three working mode of the previously described
associative memory:
\begin{description}
    \item[define\_object]: write the name of an object to the
    selected location in CAM
    \\ {\tt wa' = 0, address = name\_of\_object, sel = cam\_address
    \\ wd' = 1, din = don't\_care}

    \item[associate\_value]: write the associated value in the
    randomly accessed array to the location selected by the active
    address output of CAM
    \\ {\tt wa' = 1, address = name\_of\_object, sel = don't\_care
    \\ wd' = 0, din = value}

    \item[search]: search for the value associated with the name of the object
    applied to the address input
    \\ {\tt wa' = 1, address = name\_of\_object, sel = don't\_care
    \\ wd' = 1, din = don't\_care}
    \\ {\tt dout} is valid only if {\tt valide\_dout = 1}.
\end{description}

This associative memory will be dimensioned according to the
dimension of the actual subset of names, which is significantly
smaller than the virtual set of the possible names ($2^{p} <<<
2^{m}$). Thus, for a searching space with the size in $O(2^{m})$ a
device having the size in $O(2^{p})$ is used.






%%%%%%%%%%%%%%%%%%%%%%%%%%%%%%%%%%%%%%%%%%%%%%%%%%%%%%%%%%%%%%%%%%%%%%%%%%%%%
\section{Representation of \( B_n \) in MSB and LSB Forms with Endianness}
%%%%%%%%%%%%%%%%%%%%%%%%%%%%%%%%%%%%%%%%%%%%%%%%%%%%%%%%%%%%%%%%%%%%%%%%%%%%%
The set \( B_n \), representing \( n \)-bit unsigned binary numbers, can be expressed using two common conventions for bit ordering: Most Significant Bit (MSB) first and Least Significant Bit (LSB) first. These conventions are closely tied to the concept of endianness in digital systems, particularly when storing or transmitting binary data at the bit level. While endianness is more commonly discussed in terms of byte order in computer architecture, this section focuses on bit-level representation.

%****************************************************
\subsection{MSB Representation (Big-Endian Order)\index{Big-Endian Order}}
%****************************************************
In the Most Significant Bit (MSB)\index{MSB: most significant bit} representation, the leftmost bit (\( b_{n-1} \)) corresponds to the highest positional value (\( 2^{n-1} \)), and the rightmost bit (\( b_0 \)) corresponds to the lowest positional value (\( 2^0 \)). This notation aligns with standard positional number systems and is widely used in mathematical representations.

For \( B_n \), the value of \( x \) is computed as:
\[
x = b_{n-1} \cdot 2^{n-1} + b_{n-2} \cdot 2^{n-2} + \dots + b_1 \cdot 2^1 + b_0 \cdot 2^0, \quad b_i \in \{0, 1\}.
\]

%*******************************************************
\subsection{LSB Representation (Little-Endian Order)\index{Little-Endian Order}}
%*******************************************************
In the Least Significant Bit (LSB)\index{LSB: least significant bit} representation, the bit with the lowest weight (\( b_0 \)) is stored or transmitted first. The numerical value of the number remains the same because the bit positions still follow the standard positional notation. However, the physical storage or transmission order differs.

The mathematical formula for computing \( x \) remains:
\[
x = b_{n-1} \cdot 2^{n-1} + b_{n-2} \cdot 2^{n-2} + \dots + b_1 \cdot 2^1 + b_0 \cdot 2^0, \quad b_i \in \{0, 1\}.
\]

The distinction between MSB-first and LSB-first is primarily relevant in:
\begin{itemize}
    \item Data transmission: Some serial communication protocols transmit bits LSB-first for efficiency.
    \item Certain hardware applications: Bit-wise processing in some embedded and FPGA designs uses LSB-first order.
    \item Instruction set architectures (ISAs): Some bit manipulation instructions favor LSB-first ordering.
\end{itemize}

\textbf{Example for \( n = 3 \)}  

For \( B_3 = \{0, 1, 2, 3, 4, 5, 6, 7\} \), the logical representation (big-endian) follows:

\[
\begin{aligned}
x = 0 & \Rightarrow 000, \\
x = 1 & \Rightarrow 001, \\
x = 2 & \Rightarrow 010, \\
x = 3 & \Rightarrow 011, \\
x = 4 & \Rightarrow 100, \\
x = 5 & \Rightarrow 101, \\
x = 6 & \Rightarrow 110, \\
x = 7 & \Rightarrow 111.
\end{aligned}
\]

In a little-endian system that follows bit-order reversal, the storage or transmission order is reversed:
\begin{itemize}
    \item \( x = 1 \) (big-endian: \( 001 \)) is stored as \( 100 \) (little-endian).
    \item \( x = 6 \) (big-endian: \( 110 \)) is stored as \( 011 \) (little-endian).
\end{itemize}


\textbf{Physical Representation in Little-Endian (Stored Order)}:
\[
\begin{aligned}
x = 0 & \Rightarrow 000, \\
x = 1 & \Rightarrow 100, \\
x = 2 & \Rightarrow 010, \\
x = 3 & \Rightarrow 110, \\
x = 4 & \Rightarrow 001, \\
x = 5 & \Rightarrow 101, \\
x = 6 & \Rightarrow 011, \\
x = 7 & \Rightarrow 111.
\end{aligned}
\]



%**********************************************************************
\subsection{Mathematical Preference and Use of MSB and LSB Expansions}
%**********************************************************************

The binary expansions of \( x \) differ in the ordering of bits depending on whether the Most Significant Bit (MSB) or Least Significant Bit (LSB) form is used. These differences arise from practical applications in engineering versus conventions in mathematics.

Mathematicians universally use the MSB expansion to represent binary numbers. The MSB form is consistent with positional arithmetic, where the leftmost digit has the highest weight, and each subsequent digit corresponds to decreasing powers of the base. This notation aligns with the way numbers are taught in grade school.

In contrast, the LSB expansion prioritizes hardware efficiency, particularly in digital circuits where binary addition, shifting, and other operations naturally operate from the least significant bit upward. The storage order may change, but the arithmetic interpretation remains the same.

%*******************************************************
\subsection{Conclusion: MSB for Mathematics, LSB for Engineering}
%*******************************************************
The MSB expansion is the standard in mathematical notation due to its consistency with positional arithmetic. It provides an intuitive representation of numbers and aligns with standard conventions across different numbering systems.

The LSB-first convention, while not altering numerical value, is often employed in engineering contexts for computational efficiency, hardware design, and specific communication protocols. While endianness primarily refers to byte order in computer architecture, understanding its implications at the bit level is crucial for specialized applications such as embedded systems, networking, and low-level hardware design.

%%%%%%%%%%%%%%%%%%%%%%%%%%%%%%%%%%%%%%%%%%%%%%%%%%%%%%%%%%%%%%%%%%%%%%%%%%%%%
\section{Little-Endian vs. Big-Endian Byte Order}
%%%%%%%%%%%%%%%%%%%%%%%%%%%%%%%%%%%%%%%%%%%%%%%%%%%%%%%%%%%%%%%%%%%%%%%%%%%%%

In computer architecture, \textbf{endianness} refers to the order in which bytes are stored in memory for multi-byte data types such as integers and floating-point numbers. The two primary conventions for byte ordering are \textbf{big-endian} and \textbf{little-endian}.

%****************************************************
\subsection{Big-Endian Byte Order}
%****************************************************
In a \textbf{big-endian} system, the most significant byte (MSB) is stored at the lowest memory address, and the least significant byte (LSB) is stored at the highest memory address. This ordering aligns with how numbers are typically written in standard positional notation, where the highest-value digit appears first.

For example, consider the 32-bit hexadecimal value \( 0x12345678 \) stored in memory starting at address \texttt{0x1000} in a big-endian system:

\begin{center}
\begin{tabular}{|c|c|c|c|c|}
\hline
\textbf{Memory Address} & \texttt{0x1000} & \texttt{0x1001} & \texttt{0x1002} & \texttt{0x1003} \\
\hline
\textbf{Stored Byte}    & \texttt{0x12}   & \texttt{0x34}   & \texttt{0x56}   & \texttt{0x78}   \\
\hline
\end{tabular}
\end{center}

Big-endian ordering is used in several architectures, including older IBM mainframes, Motorola 68k processors, and certain network protocols such as TCP/IP.

%*******************************************************
\subsection{Little-Endian Byte Order}
%*******************************************************
In a \textbf{little-endian} system, the least significant byte (LSB) is stored at the lowest memory address, and the most significant byte (MSB) is stored at the highest memory address. This means that when accessing memory sequentially, the lowest-value byte appears first.

Using the same example \( 0x12345678 \), a little-endian system stores it as follows:

\begin{center}
\begin{tabular}{|c|c|c|c|c|}
\hline
\textbf{Memory Address} & \texttt{0x1000} & \texttt{0x1001} & \texttt{0x1002} & \texttt{0x1003} \\
\hline
\textbf{Stored Byte}    & \texttt{0x78}   & \texttt{0x56}   & \texttt{0x34}   & \texttt{0x12}   \\
\hline
\end{tabular}
\end{center}

Little-endian ordering is widely used in modern processors, including Intel x86, x86-64, and ARM architectures in some configurations.

%*******************************************************
\subsection{Endianness and Memory Access}
%*******************************************************
Endianness is important when interpreting multi-byte data in memory. When transferring data between systems with different endianness, byte swapping may be necessary to maintain correct values. This is particularly important in:
\begin{itemize}
    \item File formats and binary data storage: Certain file formats specify a fixed byte order, requiring conversion when read on systems with a different endianness.
    \item Networking protocols: The Internet Protocol (IP) uses big-endian byte order, often referred to as \textit{network byte order}. Little-endian systems must convert values to big-endian before transmission.
    \item Hardware and embedded systems: Some microcontrollers and peripheral devices may use a different byte ordering than the CPU, requiring explicit handling of endianness during data transfers.
\end{itemize}

%*******************************************************
\subsection{Endianness Conversion}
%*******************************************************
When transferring data between little-endian and big-endian systems, swapping byte order is necessary. In software, conversion functions such as \texttt{htonl()} (host-to-network long) and \texttt{ntohl()} (network-to-host long) are commonly used for this purpose. For example, in C:

\begin{verbatim}
uint32_t value = 0x12345678;
uint32_t swapped = __builtin_bswap32(value); // Converts between endianness
\end{verbatim}

For hardware implementations, endianness conversion is often handled using bitwise shift operations or dedicated instructions provided by the processor.

%*******************************************************
\subsection{Which Byte Endianness is Better?}
%*******************************************************
There is no inherent advantage to either endianness; the choice depends on historical and architectural considerations:
\begin{itemize}
    \item Big-endian is more intuitive when viewing hexadecimal dumps and aligns with standard numerical notation.
    \item Little-endian allows simpler hardware implementations for arithmetic operations and enables efficient incremental memory access.
\end{itemize}

Most modern systems, including x86 and ARM, use little-endian ordering due to its hardware efficiency, while big-endian is still found in network protocols and some legacy architectures.

%*************************************************
\subsection{Commentary on Bit and Byte Endianness}
%*************************************************
Big-endian systems store numbers in memory starting with the most significant bit or byte. This is often the default representation in mathematical contexts, as it aligns with traditional positional notation.

Little-endian systems store numbers starting with the least significant bit or byte. This ordering is common in modern architectures such as x86 and is often preferred for hardware operations that involve incremental processing of binary numbers, such as shift operations and carry propagation in arithmetic.

It is important to note that little-endian byte order (common in CPUs) does not reverse the bits, only the byte order. However, some hardware systems or serial communication protocols may reverse the bit order as well, depending on the application.


%%%%%%%%%%%%%%%%%%%%%%%%
\section{Problems}
%%%%%%%%%%%%%%%%%%%%%%%%
\subsection{Registers}

\begin{problem}\label{mean}
Design a system that receives a signed integer in each clock cycle and outputs the sum of the last four received numbers. When the system is reset, it is considered that the last four received numbers had the value 0.

If $t_{SU_{min}} = \#1$, $t_p = \#5$, $t_{sum} = \#50$, then what is the minimum period of the clock at which the system can present a correct result on its output.

$\diamond$
\end{problem}

{\bf Solution}

The system consists in 4 registers, {\tt reg0, ..., reg3}, serially connected and a tree of three adders used to sum the content of the registers. The first register, {\tt reg0}, receives in each clock cycle the input value, {\tt in}. Register {\tt reg1} receives the content of {\tt reg1}, {\tt reg2} receives the content of {\tt reg1}, and {\tt reg3} the content of {\tt reg2}. The first two and the last two registers are added using two adders whose outputs are added using the output adder. The division by 4 is simply performed selecting the 32 most significant bits from the output of the last adder.

\includeverilogcode[caption={Circuit that compute the mean value from the last component of a received stream of numbers.   File: A09\_mean.sv  (SystemVerilog)}, 
                    label={lst:A09_mean.sv}, 
                    firstline=1, 
                    firstnumber=1, 
                    lastline=30]
                    {\verilogSourceDir/A09_mean.sv}



\includeverilogcode[caption={Tester for circuit {\tt mean}.   File: A09\_testMean.sv  (SystemVerilog)}, 
                    label={lst:A09_testMean.sv}, 
                    firstline=1, 
                    firstnumber=1, 
                    lastline=30]
                    {\verilogSourceDir/A09_testMean.sv}


The minimal period of clock is: $T_{clock} = t_p + 2 \times t_{sum} + t_{SU_{min}} = \#106$

\begin{problem}
Revisit the Problem \ref{mean} and insert at the output of the first two adders and to the output adder pipe registers.
\begin{enumerate}
    \item What is the resulting minimal period of clock
    \item Design this pipelined version
    \item Simulate the resulting design.
\end{enumerate}

$\diamond$
\end{problem}

\begin{problem}
Design a fully buffered ALU and evaluate the minimal period of the clock signal if $t_{SU_{min}} = \#1$, $t_p = \#5$, $t_{ALU} = \#60$. Fully buffered means inputs on ALU are provided from registers and the output is send through a register.

$\diamond$
\end{problem}

\begin{problem}
Design a system that receives a string of bytes synchronized with the clock and transfers it synchronously with a selectable delay. The selection is made with {\tt sel[1:0]} indicating a delay of 1, 2, 3, or 4 clock cycles.

$\diamond$
\end{problem}

\subsection{Memories}

\begin{problem}
Design a synchronous pipelined RAM for $2^{addressSize}$ $(wordSize)$-bit words.

$\diamond$
\end{problem}

{\bf Solution}

\includeverilogcode[caption={Define the parameterized synchronous RAM.   File: A09\_defineParamSyncPipeRAM.vh  (SystemVerilog)}, 
                    label={lst:A09_defineParamSyncPipeRAM.vh}, 
                    firstline=1, 
                    firstnumber=1, 
                    lastline=30]
                    {\verilogSourceDir/A09_defineParamSyncPipeRAM.vh}

\includeverilogcode[caption={Parameterized synchronous RAM.   File: A09\_paramSyncPipeRAM.sv  (SystemVerilog)}, 
                    label={lst:A09_paramSyncPipeRAM.sv}, 
                    firstline=1, 
                    firstnumber=1, 
                    lastline=30]
                    {\verilogSourceDir/A09_paramSyncPipeRAM.sv}



\begin{problem}
Consider memory modules of 1KB and design with them a memory of 4K 32-bit word.

$\diamond$
\end{problem}

\begin{problem}
Simulate the behavior of the register file described in Section \ref{fastregFile} to prove the read-during-write operation.

$\diamond$
\end{problem} 